% =============================================================
%  1 系统网络架构
% =============================================================
\section{系统网络架构}

区块链驱动的边缘计算系统包含 $I$ 个 IoT 设备和 $J$ 个基站，
$\mathcal{I}=\{1,2,\dots,I\}$ 和 $\mathcal{J}=\{1,2,\dots,J\}$ 分别表示
物联网设备和基站的集合。该系统在逻辑上由 MEC 子系统和区块链子系统组成，
具体系统场景如图~\ref{fig:scene} 所示。

\begin{figure}[H]
  \centering
  \fbox{\parbox{0.9\linewidth}{\centering
    \vspace{2.2cm} 图~\ref{fig:scene} 占位：系统场景示意图 \vspace{2.2cm}}}
  \caption{系统场景}
  \label{fig:scene}
\end{figure}

\subsection{MEC 子系统}

在每个基站配备 MEC 服务器，可以就近为 IoT 设备提供算力。将每个基站及其所配备的
边缘服务器视为一个整体，将具备充足的存储容量和能源供应的基站作为主基站，
可以长时间保持运行状态，有效地整合和管理 IoT 设备，并提供不同层级的服务支持。

假设系统在 $\mathcal{T}$ 个时隙中运行，时隙的集合为
$\mathcal{T}=\{1,2,\dots,T\}$，在每个时隙开始时，每个 IoT 设备会生成一个任务
$W_{i}(t)=\{D_{i}(t), c_{i}(t), R_{i}(t), \tau_{i}^{\max}(t)\}$。其中
$D_{i}(t)$ 表示任务输入数据的总量，以比特为单位；$c_{i}(t)$ 表示计算复杂度，
即处理 1 比特输入数据所需的 CPU 周期数；$R_{i}(t)$ 表示处理结果的大小，
其大小远小于输入数据；$\tau_{i}^{\max}(t)$ 表示任务处理的最大容忍时延。
为确保不同任务的时延需求不会影响交叉时隙的优化问题，要求预期时间段的持续时间
大于或等于所有任务的最大允许时延。

当第 $i$ 个 IoT 设备将任务卸载给第 $j$ 个基站时，其时延包括输入数据传输时延、
基站任务处理时延以及将计算结果返回 IoT 设备的时延。考虑到任务处理结果通常小于
输入数据，其返回时延可以忽略。因此，边缘服务器处理任务的时延为

\begin{equation}
  t_{i,j}^{\mathrm{edge}}(t)
  = \frac{D_{i}(t)}{v_{i,j}(t)}
  + \frac{D_{i}(t)\, c_{i}(t)}{f_{i,j}(t)},
\end{equation}

\noindent 式中：$f_{i,j}(t)$ 表示第 $j$ 个基站分配给第 $i$ 个 IoT 设备的计算资源；
$a_{i,j}(t)$ 表示接入控制决策，若第 $i$ 个 IoT 设备接入第 $j$ 个边缘服务器，则
$a_{i,j}(t)=1$，否则 $a_{i,j}(t)=0$；$v_{i,j}(t)$ 表示第 $j$ 个基站分配给
IoT 设备的无线传输速率。

\dotsline

\subsection{区块链子系统}

每个基站既是 MEC 子系统中任务处理的边缘节点，又是区块链子系统中的区块链节点。
在区块链子系统中，主基站作为区块生成者以及共识主节点，其他基站则作为从节点
（也称为非主节点）。同时，区块生成者还可以作为区块链系统的客户，因为其生成的
区块需要经过区块链系统的共识。

采用 PBFT 协议进行区块共识包括请求、预准备、准备、提交和回复等 5 个过程，具体如下。

\subsubsection{请求}

\dotsline

\subsubsection{预准备}

\dotsline
